% Ad Familiares 4.12 (ad-familiares-04-12)
% Source: https://raw.githubusercontent.com/PerseusDL/canonical-latinLit/master/data/phi0474/phi056/phi0474.phi056.perseus-lat1.xml
% Fetched by scripts/fetch_latin.py

% Dateline (Perseus): Scr. Athenis pr. K. Iunias a. 709 (45).

\ciceroLetterOpener{SERVIVS CICERONI sal. PLVR.}

\ciceroSection{1} etsi scio non iucundissimum me nuntium vobis allaturum, tamen, quoniam casus et natura in nobis dominatur, visum est faciendum, quoquo modo res se haberet, vos certiores facere. A. d. x K. Iun. cum ab Epidauro Piraeum navi advectus essem, ibi M. Marcellum, conlegam nostrum, conveni eumque diem ibi consumpsi, ut cum eo essem. postero die cum ab eo digressus essem eo consilio, ut ab Athenis in Boeotiam irem reliquamque iuris dictionem ab solverem, ille, ut aiebat, supra Maleas in Italiam versus navigaturus erat.

\ciceroSection{2} post diem tertium eius diei cum ab Athenis proficisci in animo haberem, circiter hora decima noctis P. Postumius, familiaris eius, ad me venit et mihi nuntiavit M. Marcellum, conlegam nostrum, post cenae tempus a P. Magio Cilone, familiare eius, pugione percussum esse et duo vulnera accepisse, unum in stomacho, alterum in capite secundum aurem sperare tamen eum vivere posse; Magium se ipsum interfecisse postea; se a Marcello ad me missum esse, qui haec nuntiaret et rogaret, uti medicos ei mitterem. itaque medicos coegi et e vestigio eo sum profectus prima luce. cum non longe a Piraeo abessem, puer Acidim obviam mihi venit cum codicillis, in quibus erat scriptum paulo ante lucem Marcellum diem suum obisse.

\ciceroSection{3} ita vir clarissimus ab homine deterrimo acerbissima morte est adfectus, et, cui inimici propter dignitatem pepercerant, inventus est amicus, qui ei mortem offerret. ego tamen ad tabernaculum eius perrexi. inveni duos libertos et pauculos servos; reliquos aiebant profugisse metu perterritos, quod dominus eorum ante tabernaculum interfectus esset. coactus sum in eadem illa lectica, qua ipse delatus eram, meisque lecticariis in urbem eum referre ibique pro ea copia, quae Athenis erat, funus ei satis amplum faciendum curavi. ab Atheniensibus, locum sepulturae intra urbem ut darent, impetrare non potui, quod religione se impediri dicerent, neque tamen id antea cuiquam concesserant. quod proximum fuit, uti in quo vellemus gymnasio eum sepeliremus, nobis permiserunt. nos in nobilissimo orbi terrarum gymnasio Academiae locum delegimus ibique eum combussimus posteaque curavimus, ut eidem Athenienses in eodem loco monumentum ei marmoreum faciendum locarent. ita , quae nostra officia fuerunt pro collegio et pro propinquitate, et vivo et mortuo omnia ei praestitimus. vale . D. pr. K. Iun. Athenis.
